\documentclass[11pt,a4paper]{article}
\usepackage[T1]{fontenc}
\usepackage[utf8]{inputenc}
\usepackage[french]{babel}
\usepackage{lmodern}
\usepackage{amsmath,amssymb}
\usepackage[margin=2.6cm]{geometry}
\usepackage{booktabs}
\usepackage{enumitem}
\usepackage{xcolor}
\usepackage{hyperref}
\hypersetup{colorlinks=true,linkcolor=blue!50!black}

\newcommand{\NW}{\mathrm{NW}}
\newcommand{\E}{\mathbb{E}}
\newcommand{\code}[1]{\texttt{#1}}

\title{Résultats négatifs, artefacts et échecs de protocole\\
\large Livrable 05 du programme M3 (workspace \code{m3\_credit\_soc})}
\author{Agent de recherche M3}
\date{4 juillet 2026}

\begin{document}
\maketitle

\begin{abstract}
Tout ce qui échoue, dans le modèle comme dans le protocole. Rien ici n'est
dilué ailleurs : le rapport 04 renvoie à ce document. Les échecs sont de
trois ordres : hypothèses scientifiques non soutenues (le crédit reste
distributionnellement neutre ; Boltzmann--Gibbs échoue sur la variable même
pour laquelle il était formulé ; aucune signature SOC critique), défauts de
conception pré-enregistrée (l'ablation F éteint le marché au lieu de tester
la topologie ; G1 confond corrélation et raréfaction du crédit ; le canal
« défaut de liquidité », pièce maîtresse de M3, ne s'active dans aucun
régime testé), et artefacts d'implémentation détectés puis corrigés (bug
d'observation I8).
\end{abstract}

\section{Hypothèses scientifiques non soutenues}

\subsection{Le crédit productif reste distributionnellement neutre
(l'objectif central de M3 n'est pas atteint)}
M3 a été construit pour donner au crédit un canal causal réel : conversion
$L \to K$ exclusive, obligations exigibles en liquidité, actif illiquide.
Résultat (10 seeds A vs 10 seeds B, mêmes seeds, $T{=}4000$) : les formes
sont indiscernables — corps de $\NW$ Fisk des deux côtés, exposant de queue
$2{,}69$ des deux côtés (égalité au centième), revenu dPlN 10/10 des deux
côtés, $K$ lognormal, Gini et parts du top égaux. Le crédit modifie des
\emph{niveaux démographiques} (population et espérance de vie $\times 0{,}5$,
renouvellement accéléré) mais aucune \emph{forme}. La conclusion négative de
M2 (« le crédit n'a pas d'effet distributionnel ») survit donc à un modèle
conçu explicitement pour la renverser. Nuance importante : le critère causal
pré-enregistré (deux signatures modifiées) est \emph{formellement} atteint —
mais par des signatures redondantes entre elles (toutes des facettes de
« vies plus courtes »). Si l'on avait voulu que ce critère mesure un effet
distributionnel, il fallait l'écrire autrement : leçon de pré-enregistrement.

\subsection{Boltzmann--Gibbs échoue sur $L$, sa propre variable}
M2 avait testé l'exponentielle sur $\NW$ — transposition abusive dénoncée
dans le rapport 01. M3 a introduit $L$ précisément pour poser la question
correctement. Verdict : corps de $L$ Fisk (expon perd de $\sim 700$ points
d'AIC ; $\mathrm{med/mean} = 0{,}93$ au lieu de $\ln 2 \approx 0{,}69$).
Le diagnostic est structurel, pas paramétrique : l'injection $(1-s)P_i$ est
hétérogène (proportionnelle à $\sqrt{K_i}$) et la fuite $cL$ multiplicative ;
les échanges conservatifs (intérêts, $1{,}5\,\%$ de la production) sont
négligeables devant ces flux. Dans toute cette classe de modèles
production--consommation, l'hypothèse de conservation locale est violée par
construction : le cadre YR n'y a \emph{pas de domaine d'application}. Tester
« l'exponentielle de la monnaie » dans un modèle où la monnaie est
essentiellement un flux de production recyclé n'a pas de sens — à retenir
pour M4.

\subsection{Le canal « défaut de liquidité » ne s'active jamais}
La fragilité nominal/réel (défaut d'une entité solvable mais illiquide) est
le mécanisme différenciant de M3. Il est implémenté, testé unitairement
(50 tests), mesurable (comptage des racines par cause) — et \textbf{ne se
déclenche dans aucun run} : 0 défaut de liquidité sur $\sim 200\,000$
pas-entité$\times 10^3$ cumulés (baseline, B--J, chocs corrélés compris).
Cause mesurée : la charge d'intérêts d'équilibre est trop faible
($1{,}5\,\%$ de la production ; les emprunteuses convergent vite vers
$K^*(\rho) \approx 23$ et leur afflux de liquidité couvre toujours le
service). Verdict de la grille complète (45 cellules) : 4 événements en
tout, uniquement à $s = 0{,}9$ (afflux de liquidité minimal), zéro partout
ailleurs, chocs corrélés compris. Seule exception notable : l'ablation H
($d_0 = 0$), où le levier relatif est bien plus élevé, en produit 71 —
c'est un régime différent, pas un réglage de la baseline. Conclusion de
conception pour M4 : un marché du crédit \emph{au taux marginal
concurrentiel}, servi par une liquidité alimentée par la production
courante, ne produit pas d'illiquidité ; il faudrait des taux
hors-équilibre, du principal exigible (dette amortissable, pas
perpétuelle), des appels de marge — ou le plancher endogène de H.

\subsection{Aucune signature SOC critique}
Avalanches causales (définition corrigée : composantes du graphe de pertes,
pas les lots par pas de M2) : sous-Poisson en baseline (var/moy $0{,}06$,
max 7), premières rafales sous chocs sectoriels G2 (max 36, var/moy
$0{,}38$) mais toujours sous-critiques (singletons $> 98\,\%$, var/moy
$< 1$). corr(concentration $\to$ faillites futures) systématiquement
\emph{négative} ($-0{,}2$ à $-0{,}6$) : le signe même de
l'accumulation--relaxation est absent — la concentration du crédit suit les
périodes calmes, elle ne les précède pas. Le pouvoir prédictif individuel
du réseau est réel mais modeste (AUC levier $\to$ mort $0{,}59$--$0{,}64$).

\subsection{Sans renouvellement, extinction (ablation I)}
$\lambda = 0$, cohorte initiale 1000 : extinction totale à $t \approx 1760$
(2 seeds sur 3 ; le troisième à $N = 1$ à $T = 2000$). Le mécanisme
démographique Reed n'est pas une composante de la distribution : c'est la
condition d'existence du système. Corollaire : aucune « distribution
stationnaire de M3 » n'existe indépendamment du flux de naissances.

\section{Échecs de protocole pré-enregistré}

\subsection{L'ablation F ne teste pas ce qu'elle devait tester}
F (appariement entièrement aléatoire) devait isoler la topologie à bilans
agrégés comparables. Mesure : volume $11{,}6$/pas contre $75{,}9$ en
baseline, 452 contrats actifs contre $68\,608$ — la condition « bilans
comparables » est massivement violée ; F reproduit B en éteignant le marché.
La question « la topologie compte-t-elle ? » reste \emph{sans réponse
confirmatoire} dans ce programme. Une variante F$''$ (emprunteuse
assortative, prêteuse aléatoire parmi les faisables) est exécutée en
exploratoire, clairement étiquetée post-hoc.

\subsection{G1 confond corrélation macro et raréfaction du crédit}
Le design pré-enregistré maintient $\sigma^2$ total constant ; la part macro
réduit donc mécaniquement la dispersion idiosyncratique des $K$, qui est le
moteur du gain à l'échange — le crédit s'étiole (656--3676 contrats contre
$68\,608$). L'effet « chocs corrélés » et l'effet « moins de crédit » sont
inséparables dans G1. G2 (sectoriel) souffre moins de ce biais (crédit
intact) et porte seul le verdict sur la corrélation.

\subsection{Population divergente hors de la colonne $c = 0{,}10$}
La fenêtre démographique de M3 est plus étroite que prévu : sur les plages
pré-enregistrées $s \in [0{,}5, 0{,}9] \times c \in [0{,}02, 0{,}10]$, seule
la colonne $c = 0{,}10$ est viable (8 cellules sur 12 divergent,
morts/naissances $0{,}34$--$0{,}72$). Tout stock stable ajouté à $\NW$
amortit le plancher $d_0$ et éteint la mortalité — la même fragilité de
fenêtre que M2 ($\sigma$), déplacée sur $c$. Le « sans réglage fin » reste
donc partiel : la \emph{forme} est robuste, l'\emph{existence du régime}
dépend d'un paramètre à un facteur 2 près.

\section{Artefacts d'implémentation détectés et corrigés}

\subsection{Observation non neutre (violation de I8)}
La fonction d'instantané lisait les agrégats du carnet via \code{[]} sur des
\code{defaultdict}, insérant des clés vides et modifiant l'ordre d'itération
du service des intérêts : un run \emph{avec} snapshots divergeait au dernier
ulp d'un run \emph{sans} (détecté par le test I8 strict, seed 9, dès
$t = 54$). Corrigé (accès en lecture pure) ; I8 bit à bit rétabli. Leçon
générique : sur un \code{defaultdict}, toute lecture par \code{[]} est une
écriture ; les fonctions d'observation doivent être prouvées neutres, pas
supposées telles.

\subsection{Estimateur CSN dégénéré sur les tailles d'avalanches}
\code{fit\_tail\_csn} divisait par zéro quand toutes les valeurs de queue
égalent $x_{\min}$ (cas réel : ablation B, toutes les avalanches de taille
1). Garde ajoutée (candidat écarté si le dénominateur log est nul). Par
ailleurs les « ajustements de queue » sur des tailles d'avalanches presque
toutes égales à 1 produisent des exposants absurdes ($\alpha \approx 50$) :
ils ne sont \emph{pas} interprétés dans les rapports — seuls max, var/moy et
fraction multi le sont.

\subsection{Fragilité du sous-agent d'exécution}
Anecdote de méthode, consignée pour la reproductibilité du flux de travail :
la délégation de la surveillance des runs à un agent léger a échoué deux
fois (l'agent « attendait » sans commande bloquante) ; les runs eux-mêmes,
lancés par lui, étaient corrects. La surveillance a été reprise par le
processus principal sans perte de calcul.

\section{Ce que ces négatifs enseignent}
Le noyau Reed (chocs multiplicatifs, plancher, réinjection) reste, comme en
M2, le seul générateur des formes observées ; un crédit productif au taux
marginal, servi en liquidité abondante, n'y ajoute qu'une accélération
démographique. Pour qu'un M4 donne sa chance à la thèse
accumulation--relaxation, il faudra au moins l'un de : dette amortissable
(principal exigible, pas perpétuel), taux hors-équilibre ou spread de
risque endogène, liquidité réellement rare (injection découplée de la
production courante), ou levier avec appel de marge. Sans quoi le résultat
« le crédit ne fait que raccourcir les vies » se répétera.

\end{document}
